\SourceLeaf{269}%
auf den Abstand vom mittlern Parallelkreise, und enthält ausserdem noch die Abplattung oder das Quadrat von \(e\) als Factor.

Allein dieser Vortheil lässt sich noch sehr vergrössern, wenn man anstatt jener Bestimmung der willkürlichen Function eine etwas abgeänderte, für die Rechnung fast eben so bequeme wählt, indem man nemlich unter Zuziehung einer zweiten Constante \(\alpha\),
\[
f\upsilon=\alpha\upsilon-i\log k
\]
setzt. Man hat dann in seiner Gewalt, durch zweckmässige Bestimmung der beiden Constanten zu bewirken, dass die Abweichung des Vergrösserungsverhältnisses \(m\) von dem Werthe 1, in Beziehung auf den Abstand vom mittlern Parallelkreise eine Grösse der dritten Ordnung wird, ungerechnet den auch hier bleibenden Factor \(ee\).

\Article{3}
\indent Die Formel 1 gibt, bei dieser Bestimmung der Function \(f\).
\[
T=\alpha t \quad .\ .\ .\ .\ .\ .\ .\ .\ .\ .\ .\ .\ .\ .\ .\ .\ .\quad(2)
\]
\[
\operatorname{tang}\tfrac{1}{2}U=k\operatorname{tang}\tfrac{1}{2}w^\alpha(\tfrac{1+e\cos w}{1-e\cos w})^{\frac{1}{2}\alpha e} \quad .\ .\ .\ .\ .\ .\ .\ .\ .\quad(3)
\]
und für \(m\) findet man leicht, aus den in der mehrerwähnten Schrift entwickelten Grundsätzen, den Ausdruck
\[
m=\frac{\alpha A\sin U\BareRoot{(1-ee\cos w^2)}}{a\sin w} \quad .\ .\ .\ .\ .\ .\ .\ .\ .\ .\quad(4)
\]
wenn durch \(a\) die halbe grosse Achse des Ellipsoids und durch \(A\) der Halbmesser der Kugel bezeichnet wird.

Die Differentiation der logarithmisch ausgedrückten Gleichung 3 ergibt
\[
\frac{\mathrm{d}U}{\sin U}=\frac{\alpha\mathrm{d}w}{\sin w}-\frac{\alpha ee\sin w\mathrm{d}w}{1-ee\cos w^2}
\]
oder
\[
\frac{\mathrm{d}U}{\mathrm{d}w}=\frac{\alpha(1-ee)\sin U}{(1-ee\cos w^2)\sin w} \quad .\ .\ .\ .\ .\ .\ .\ .\ .\ .\ .\ .\quad(5)
\]
Ebenso ergibt die Differentiation der Gleichung 4
\[
\begin{aligned}
\mathrm{d}\log m&=\operatorname{cotg}U\mathrm{d}U-\operatorname{cotg}w\mathrm{d}w+\frac{ee\cos w\mathord{.}\sin w\mathrm{d}w}{1-ee\cos w^2}\\
&=\operatorname{cotg}U\mathrm{d}U-\frac{(1-ee)\cos w\mathrm{d}w}{(1-ee\cos w^2)\sin w}
\end{aligned}
\]
